% !PRINTED 557
\weberpage{557}

%<<CONT
(10) is isomorphic with the permutation group $C$. Hence $C$ is the \gesperrt{Galois} group of $X = 0$.

By \S.\,163 the equation $X = 0$ can be reduced to a sequence of cyclic equations of lower degree, if one seeks functions that belong to the divisors of the group $C$. But the divisors of $C$, if $n - 1$ is decomposed into two factors $e, f$, so that
\[
n - 1 = ef
\]
is set, are the cyclic groups

\weq{11}{C_e = 1, \pi^e, \pi^{2e} \ldots \pi^{(f-1)e},}

and the permutation $\pi^e$ is therein equivalent to the substitution $(r, r_e)$.

In order to arrive at a function of the roots belonging to $C_e$ of the simplest possible form, we consider, following \gesperrt{Gauss}, the $f$-\gesperrt{term periods}

\weq{12}{\begin{aligned}
\eta &= r + r_e + r_{2e} + \cdots + r_{(f-1)e} \\
\eta_1 &= r_1 + r_{e+1} + r_{2e+1} + \cdots + r_{(f-1)e+1} \\
&\cdots\cdots\cdots\cdots\cdots\cdots\cdots\cdots \\
\eta_{e-1} &= r_{e-1} + r_{2e-1} + r_{3e-1} + \cdots + r_{n-2},
\end{aligned}}

which arise from the first of them by the substitutions $(r, r), (r, r_1), (r, r_2) \ldots (r, r_{e-1})$. We also designate the quantities (12) as a system of \gesperrt{conjugate periods}.

Each of these functions remains unchanged by the substitution $(r, r_e)$, and they are all different from one another, since an equation $\eta_h - \eta_k = 0$, if $h$ is different from $k$, could be written in the form
\[
a r + a_1 r_1 + \cdots + a_{n-2} r_{n-2} = 0
\]
where the $a, a_1 \ldots a_{n-2}$ are integers not all vanishing at once ($+ 1$, $0$ or $- 1$), which by Theorem I is not possible. Each of the periods (12) is therefore a function \gesperrt{belonging} to the group $C_e$, and these quantities are therefore the roots of an irreducible equation of degree $e$ with integer coefficients. If one adjoins one of these quantities, the determination of $r$ still depends on an equation of degree $f$.

By the general theorems of the fourteenth section, every function of the $r$ that admits the substitutions of the group $C_e$ can be represented rationally by any one of the quantities $\eta$. But here we can also set up the following theorem:
%CONT>>
